\documentclass{nbeamer}
\usepackage{nbeamer-code}

\title{nbeamer 是一个极简风的 \LaTeX{} beamer class}
\subtitle{nbeamer 使用手册}
\author{XIAN}
\institute{纳米树研究室 Namishu Research Lab}
\date{\today}

\begin{document}
\begin{frame}
  \titlepage
\end{frame}

\begin{frame}
  \title{目录}
  \usebeamertemplate{contents page}
\end{frame}

\section{开始}
\begin{frame}[fragile]{最小示例}
保存为 \incode{talk.tex}。正文页只显示主标题。使用 \incode{\framesubtitle} 会收到编译警告。
\begin{ncode}
\documentclass{nbeamer}
\begin{document}
\begin{frame}{第一页}
  你好，nbeamer。
\end{frame}%
\end{document}
\end{ncode}
在同一目录运行 \incode{latexmk -xelatex talk.tex}，打开 \incode{talk.pdf}。
\end{frame}

\section{页面结构}
\begin{frame}[fragile]{封面}
本手册首页的五个字段写在导言区。\incode{\today} 显示年-月-日，也可在 \incode{\date{...}} 中自填日期。
\begin{ncode}
\title{nbeamer 是一个极简风的 \LaTeX{} beamer class}
\subtitle{nbeamer 使用手册}
\author{XIAN}
\institute{纳米树研究室 Namishu Research Lab}
\date{\today}
\end{ncode}
在正文中用一个 frame 放置封面：\incode{\begin{frame}\titlepage\end{frame}}。
\end{frame}

\begin{frame}[fragile]{目录}
目录收集各个 \incode{\section} 的名称。在封面后插入目录，在各组正文前用这个命令命名章节。
\begin{ncode}
\begin{frame}
  \title{目录}
  \usebeamertemplate{contents page}
\end{frame}%

\section{正文元素}
\end{ncode}
\end{frame}

\section{正文元素}

\begin{frame}[fragile]{文字强调}
\begin{table}
  \begin{tabular*}{\linewidth}{l@{\extracolsep{\fill}}ll}
    \toprule
    来源 & 写法 & 效果 \\
    \midrule
    LaTeX   & \incode{\textbf{关键词}}       & \textbf{关键词} \\
            & \incode{\emph{关键条件}}        & \emph{关键条件} \\
    Beamer  & \incode{\alert{仅限首次提交}}   & \alert{仅限首次提交} \\
    nbeamer & \incode{\nbemph{主要结论}}      & \nbemph{主要结论} \\
            & \incode{\nbmute{数据来源}}      & \nbmute{数据来源} \\
    \bottomrule
  \end{tabular*}
\end{table}

\incode{\alert} 是 Beamer 自带的提醒命令，nbeamer 将其样式设为砖红色粗体。
\end{frame}

\begin{frame}[fragile]{列表}
\begin{columns}
\begin{column}{13\nbgrid}
\begin{itemize}
  \item 无序条目
  \begin{itemize}
    \item 子项
  \end{itemize}
\end{itemize}
\begin{ncode}
\begin{itemize}
  \item 无序条目
  \begin{itemize}
    \item 子项
  \end{itemize}
\end{itemize}
\end{ncode}
\end{column}
\begin{column}{13\nbgrid}
\begin{enumerate}
  \item 第一项
  \item 第二项
\end{enumerate}
\begin{ncode}
\begin{enumerate}
  \item 第一项
  \item 第二项
\end{enumerate}
\end{ncode}
\end{column}
\end{columns}
\end{frame}

\begin{frame}[fragile]{内容块}
\begin{block}{命题}
两个偶数之和是偶数。
\end{block}
\begin{proof}[证明]
两个加数都能被 2 整除，它们的和也能被 2 整除。
\end{proof}
\begin{ncode}
\begin{block}{命题}
  两个偶数之和是偶数。
\end{block}
\begin{proof}[证明]
  两个加数都能被 2 整除，它们的和也能被 2 整除。
\end{proof}
\end{ncode}
\end{frame}

\begin{frame}[fragile]{公式}
行内公式跟随文字，例如 $a^2+b^2=c^2$。
\begin{ncode}
$a^2+b^2=c^2$
\end{ncode}

同一个公式也可以单独居中：
\[
  a^2+b^2=c^2
\]
\begin{ncode}
\[
  a^2+b^2=c^2
\]
\end{ncode}
\end{frame}

\begin{frame}[fragile]{图片}
\begin{columns}
\begin{column}{11\nbgrid}
\begin{figure}
  \includegraphics
    [width=\linewidth]
    {figures/mushroom-shelter.png}
  \caption{蘑菇下的两只鸟}
\end{figure}
\end{column}
\hfill
\begin{column}{15\nbgrid}
\begin{ncode}
\begin{figure}
  \includegraphics
    [width=\linewidth]
    {figures/mushroom-shelter.png}
  \caption{蘑菇下的两只鸟}
\end{figure}
\end{ncode}
\end{column}
\end{columns}
\end{frame}

\begin{frame}[fragile]{表格}
\begin{columns}
\begin{column}{8\nbgrid}
\begin{table}
  \begin{tabular*}{\linewidth}
    {l@{\extracolsep{\fill}}c}
    \toprule 格式 & 对齐方式 \\
    \midrule
    l & 左对齐 \\
    c & 居中 \\
    r & 右对齐 \\ \bottomrule
  \end{tabular*}
  \caption{列格式}
\end{table}
\end{column}
\hfill
\begin{column}{15\nbgrid}
\begin{ncode}
\begin{table}
  \begin{tabular*}{\linewidth}
    {l@{\extracolsep{\fill}}c}
    \toprule 格式 & 对齐方式 \\
    \midrule
    l & 左对齐 \\
    c & 居中 \\
    r & 右对齐 \\ \bottomrule
  \end{tabular*}
  \caption{列格式}
\end{table}
\end{ncode}
\end{column}
\end{columns}
\end{frame}

\section{nbeamer 命令}

\begin{frame}[fragile,nbeamer={grid=true,boxes=true}]{排版辅助线}
\begin{ncode}
\nbsetup{grid=true,boxes=true}
\end{ncode}
在导言区写入这条命令，即可为整份文稿开启排版辅助线。
\incode{grid} 显示网格，方便数出栏宽和留白各占几格；
\incode{boxes} 标出标题区和内容区的边界。

\begin{ncode}
\begin{frame}[nbeamer={grid=true,boxes=true}]{分栏}
\end{ncode}
写在 \incode{frame} 的选项中，设置只对这一页生效。
\end{frame}

\begin{frame}[fragile]{网格间距}
\incode{\nbgrid} 表示一格的长度，等于幻灯片宽度的 $1/32$。

整页宽 32 格。默认左右页边距各占 2 格，中间的内容区宽 28 格。

\begin{itemize}
  \item \incode{\vspace{\nbgrid}} 放在两段之间，额外留出一格竖直间距。
  \item \incode{\hspace{2\nbgrid}} 放在同一行的文字之间，留出两格水平间距。
  \item \incode{\\[2\nbgrid]} 在段内换行，并在正常行距上额外留出两格。
  \item \incode{\nbgrid} 也可以指定栏宽，例如 \incode{\begin{column}{13\nbgrid}}。
\end{itemize}
\end{frame}

\end{document}
